% 表内破折号只表示模板尚未填入数据，不表示零、未检出或缺失观测。
% 定稿时应在表注中明确定义实际采用的缺失值/未检出标记。
\chapter{研究结果与讨论}
\section{运行表现与处理效果}
本节按研究目标呈现主要结果，说明数据范围、运行阶段及评价依据。图表应在正文首次提及后编排，内容与正文相互补充。主要指标可按表\ref{tab:results}组织。
\begin{table}[htbp]
\caption{主要指标记录表}\label{tab:results}
\centering\tjsmall
\begin{tabularx}{\textwidth}{@{}lXXX@{}}
\toprule
指标 & 进水 & 出水 & 评价说明\\\midrule
COD / (\unit{\milli\gram\per\litre}) & — & — & —\\
TN / (\unit{\milli\gram\per\litre}) & — & — & —\\
\bottomrule
\end{tabularx}
\end{table}
\section{影响因素与机制讨论}
本节结合对照试验、物质平衡、表征或模型结果解释主要现象，讨论影响因素及可能机制。研究装置或处理流程可按图\ref{fig:route}组织。
\begin{figure}[htbp]
\centering
% 正式图片写法：\includegraphics[width=.85\textwidth]{figures/文件名.pdf}
\fbox{\parbox[c][28mm][c]{.85\textwidth}{\centering\normalsize 研究装置或工艺流程图}}
\caption{工艺流程图}\label{fig:route}
\end{figure}
\section{本章小结}
本节归纳本章主要发现，说明结论的适用条件及其与后续工程评价的联系。
